%!TEX root = ../template.tex
%% ====================================================================
%% Apendice B -- Extracao assistida e validacao. Escrito em 2026-09-17.
%% Fontes: LocalResearch/Screening/level3_extraction (VALIDATION_PROTOCOL_C.md,
%% STAGE_C_2026-09-12.md, PILOT_RESULTS_2026-09-12-1.md, CORPUS_CLOSURE_2026-09-12.md,
%% production/reviewer_workspace/reports/TIER1_CLOSURE_2026-09-09.md,
%% stage_c/QUARANTINE_472536714_2026-09-15.md, stage_c/verification_v1/).
%% Tabela B.3 gerada por stage_c/figuras_cap3/tabela_apendice_b.py.
%% Versao anterior (esqueleto) em _revisao/2026-09-17-apB-antes.tex.
%% ====================================================================

\chapter{Extração assistida e validação}
\label{app:validacao}

Este apêndice descreve como foram extraídos os dados da revisão de âmbito do
Capítulo~\ref{cha:revisao}, que controlos se aplicaram e o que cada controlo
permite afirmar. A regra que o organiza é a mesma do capítulo: um número vale
o que vale a camada de onde vem. Os desvios ao protocolo registado estão no
Apêndice~\ref{app:protocolo-sr}; aqui descreve-se o que foi feito.

\section{Três camadas}
\label{sec:apB-camadas}

A extração decorreu em três camadas, com instrumentos e níveis de decisão
humana diferentes (Tabela~\ref{tab:apB-camadas}).

\begin{table}[htbp]
  \centering\footnotesize
  \caption[Camadas da extração e o que cada uma permite afirmar]{Camadas da
  extração, instrumento, decisão humana e alcance das afirmações.}
  \label{tab:apB-camadas}
  \begin{tabularx}{\textwidth}{@{}>{\raggedright\arraybackslash}p{2.2cm}>{\raggedright\arraybackslash}X>{\raggedright\arraybackslash}p{3.0cm}>{\raggedright\arraybackslash}X@{}}
    \toprule
    Camada & Instrumento & Decisão humana & Alcance \\
    \midrule
    Pertença aos corpora & Pré-extração por duas famílias de modelos; formulário às cegas para as divergências & Por publicação, nas 82 com divergência & Contagens de A e de B, e a pertença de cada publicação \\
    \addlinespace
    Campos do Corpus~A & Duas famílias de modelos, esquema v2.2.1; só valores concordantes & Nenhuma por valor; auditoria por amostragem com estado por confirmar & Distribuições descritivas do subconjunto concordante, com denominador \\
    \addlinespace
    Campos do Corpus~B & Uma família de modelos, esquema v2.3.0-C, citação e página por valor; releitura por modelo em 2026-09-17 & Nos itens que a regra de adjudicação não resolveu; amostra de controlo por rever & Contagens e ordens do Corpus~B, com as ressalvas da Secção~\ref{sec:apB-verificacao} \\
    \bottomrule
  \end{tabularx}
\end{table}

\section{Superfícies e proveniência}
\label{sec:apB-superficies}

Os modelos correram em aplicações de subscrição e não por API: o Claude Code,
com o modelo reportado em cada execução como \texttt{claude-opus-4-8} (e um
modelo auxiliar \texttt{claude-haiku-4-5}), e o Codex CLI numa subscrição
ChatGPT Plus. Estas superfícies não expõem parâmetros de amostragem, versões
datadas dos modelos nem identificadores de pedido, pelo que a proveniência se
regista ao nível da superfície: versão da aplicação, modelo reportado, data e
\emph{hash} do instrumento. Cada publicação foi tratada numa sessão nova, sem
ferramentas, a partir do texto extraído do PDF página a página. Uma
reexecução com versões futuras dos modelos não reproduz necessariamente os
valores; o que fica reprodutível é o instrumento, os registos e as
verificações feitas sobre eles. O revisor é o único adjudicador, e os
modelos são instrumentos de extração, não revisores independentes: dois
modelos podem partilhar o mesmo erro, e a concordância entre eles é usada
para encaminhar casos, nunca como medida de fiabilidade.

\section{Pertença aos corpora}
\label{sec:apB-pertenca}

A pertença aos corpora foi pré-extraída pelas duas famílias. Nas 82
publicações em que as extrações divergiam num campo com consequência para a
pertença (90 células), o revisor decidiu num formulário às cegas, sem ver os
valores propostos. O protocolo previa que, depois de gravada a decisão cega,
o formulário revelasse a proposta dos modelos e registasse qualquer
alteração. Esse passo foi executado integralmente numa publicação, sem
alterações, e parcialmente em três; em cinco a proposta foi revelada por uma
sessão automatizada, que alterou um valor, revertido a partir da cópia cega
imutável; nas restantes 73 não foi executado, e a decisão cega foi assumida
como final por decisão do revisor de 2026-09-09. Nenhuma decisão de pertença
foi tomada com a proposta dos modelos à vista; perdeu-se a segunda leitura
sobre 81 das 82 publicações e, com ela, a medida de quanto essa leitura teria
mudado as decisões.

O controlo de falsos negativos do Corpus~B e os sete veredictos adjudicados
sobre excertos literais, sem abertura dos PDF completos, estão descritos no
desvio D6 (Apêndice~\ref{app:protocolo-sr}); nenhum falso negativo foi
confirmado. A correção de pertença de 2026-09-17, que retirou do Corpus~B dois
relatórios excluídos de A por falta de modelo, está na
Secção~\ref{sec:pesquisa}; a regra de que B é um subconjunto de A está
imposta no código que gera os corpora.

\section{Campos do Corpus~A}
\label{sec:apB-corpus-a}

Os campos do Corpus~A (tipo de estudo, setor, tipo de organização,
enquadramento normativo, protocolo de medição e verificação, famílias de
modelos, país) foram extraídos pelas duas famílias com o esquema v2.2.1,
congelado antes da extração. A análise usa só os valores concordantes; as
divergências, as extrações únicas e as ausências de valor ficam como ``por
resolver'' e aparecem assim nas figuras. A cobertura por campo está na
Secção~\ref{sec:corpus-a}: nas famílias de modelos, o campo mais próximo da
SQ1, há 162 valores concordantes em 331.

Cada valor devia trazer uma citação do texto. A verificação de ancoragem
desta camada aceitava citações por semelhança, com limiares de 95, 90 e 85
consoante o comprimento, e 28\% dos valores não tinham citação e ficaram
fora dessa garantia. A auditoria por amostragem aleatória simples prevista na
emenda C-2.3 (40 publicações, semente 44), que daria um intervalo de erro por
campo, tem o estado por confirmar (caixa AP-A-04 no
Apêndice~\ref{app:protocolo-sr}). Sem ela, não há limite de erro para os
valores concordantes do Corpus~A.

\section{Campos do Corpus~B: o instrumento Stage~C}
\label{sec:apB-stage-c}

\paragraph{Instrumento.} O Corpus~B foi extraído de novo com o esquema
v2.3.0-C, congelado a 2026-09-13 com o \emph{hash} registado na emenda
(Apêndice~\ref{app:protocolo-sr}). O esquema organiza cada registo em
publicação, aplicações e modelos, para que os valores de dois modelos do
mesmo artigo não se misturem na mesma linha. Os campos C1 a C8 descrevem cada
modelo (variável dependente, forma, variáveis explicativas, dados, medição na
fronteira, regime de re-estimação, domínio de validade e métricas); C9 e C11
codificam mecanismos e adaptações; C10, C12 e C13 registam limitações,
requisitos e reprodutibilidade. C9 e C11 foram derivados do pré-preenchimento
das 42 publicações então no Corpus~B, antes do congelamento: é uma revisão do
instrumento informada pelo corpus, declarada como tal. O mesmo texto de
instruções serviu as duas superfícies.

\paragraph{Ancoragem.} Todos os valores têm uma citação literal e a página
física. Distinguem-se três verificações: a citação existe; está na página
indicada; e sustenta o valor. As duas primeiras são automáticas e exatas,
depois de duas regras de normalização publicadas (hifenização e ligaduras;
colunas intercaladas e números de linha dos manuscritos). O localizador de
página pode ser corrigido quando a citação existe palavra por palavra numa
única outra página (desvio D5). A terceira verificação não é automática e não
se afirma para esta camada senão através da releitura descrita na
Secção~\ref{sec:apB-verificacao}.

\paragraph{Piloto.} Antes do congelamento, cinco publicações heterogéneas
(refinação, petroquímica, vidro e ferro fundido, nove setores, cimento)
foram extraídas com o esquema provisório. As cinco validaram; as 538
citações foram encontradas na página indicada, e as 184 ausências declaradas
tinham base de pesquisa registada. Doze das dezassete tentativas foram gastas
em cinco defeitos do instrumento --- regras que nenhuma extração podia
satisfazer ---, todos corrigidos com um teste próprio antes do congelamento.
Os registos do piloto não são registos de extração: as cinco publicações
foram extraídas de novo depois do congelamento.

\paragraph{Produção e quarentena.} A extração correu numa só superfície
(desvio D4), uma sessão por publicação. Das 42 publicações então no
Corpus~B, 41 validaram. A publicação \texttt{rayyan-472536714}, um estudo de
\emph{benchmarking} com pelo menos quinze aplicações, não foi devolvida
inteira em quatro tentativas: a resposta vinha truncada no início, sem as
primeiras treze aplicações. Um analisador mais tolerante recuperaria um
fragmento válido e registá-lo-ia como a extração do artigo, com treze setores
a menos; por isso ficou em quarentena, e a razão é uma limitação da
superfície, não do conteúdo. Depois da correção de pertença, o Corpus~B tem
40 publicações, das quais 39 com registo.

\section{Verificação por modelo}
\label{sec:apB-verificacao}

A verificação humana por ocorrência prevista para o Corpus~B não foi feita
(desvio D7). Em sua substituição, a 2026-09-17, as células com citação de que
os resultados da Secção~\ref{sec:corpus-b} dependem foram relidas por um
segundo passe de modelo.

\paragraph{Âmbito e procedimento.} Entraram 368 células: 145 mecanismos, 91
adaptações e 132 valores de modelo (fronteira, re-estimação, domínio,
validação fora da amostra e reprodutibilidade). Foram distribuídas por sete
sessões independentes de um modelo da mesma família do extrator, cada uma
com a página citada e as adjacentes e com um livro de códigos escrito a partir
das instruções de extração. Para cada célula, o verificador respondia a três
perguntas: a citação está na página; lida no seu parágrafo, sustenta o valor;
se não sustenta, qual é o valor certo, escolhido só entre os permitidos. O
livro de códigos pedia uma decisão conservadora --- não inferir causas que o
artigo não nomeia, marcar dúvida sempre que a confiança não fosse alta ---,
porque uma confirmação errada passaria sem revisão.

\paragraph{Regra de adjudicação.} As propostas foram triadas por uma regra
fixada antes de as ler: aceitam-se as correções que tornam o valor mais fraco
ou que aplicam uma regra explícita do livro de códigos; as que o tornam mais
forte só com evidência quantitativa na página ou citação da norma; o resto
fica para o revisor.

\paragraph{Resultados.} O verificador confirmou 225 valores com confiança
alta e 7 com confiança baixa, propôs correção em 104 e remoção em 21, e deu 11
como não sustentados sem valor alternativo. Aplicada a regra, o estado final
à data desta versão é o da Tabela~\ref{tab:apB-verificacao}.

\begin{table}[htbp]
  \centering\small
  \caption[Estado das células do Corpus B após a verificação por
  modelo]{Estado das 368 células do Corpus~B após a verificação por modelo e
  a adjudicação de 2026-09-17. Pendentes: itens de dúvida D1--D8 e regras R1
  e R4, que ficam no valor original até decisão do revisor.}
  \label{tab:apB-verificacao}
  \begin{tabular}{lrrrrrr}
    \toprule
    Tipo & Confirmado & Corrigido & Retirado & Mantido & Pendente & Total \\
    \midrule
    Mecanismos  & 87  & 34 & 13 & 5 & 6  & 145 \\
    Adaptações  & 49  & 24 & 3  & 0 & 15 & 91 \\
    Modelos     & 89  & 33 & 0  & 0 & 10 & 132 \\
    \midrule
    Total       & 225 & 91 & 16 & 5 & 31 & 368 \\
    \bottomrule
  \end{tabular}
\end{table}

\paragraph{O que fica para o revisor.} Os treze itens de dúvida agrupam-se em
oito perguntas: se os pressupostos de um modelo de simulação são mecanismos
(D1); que efeito atribuir a um consumo desligado da produção (D2); se a única
menção a incrustação é um mecanismo ou uma motivação (D3); se os modelos
ajustados por modo e crude de um caso de fornecedor são re-estimação em linha
ou referência segmentada (D4); se metas ajustadas pelos últimos três meses
são uma janela móvel da referência (D5); se uma função definida para
produção positiva declara um domínio de validade (D6); a que modelos se
aplica uma frase sobre limites de validade (D7); e se dados e código em
apêndice contam como disponíveis (D8). As duas regras pendentes são o
critério de enquadramento normativo implícito das adaptações (R1, catorze
itens) e se o consumo não produtivo encontrado nos dados é um mecanismo (R4,
dois itens). Uma amostra aleatória de dez confirmações (semente 44) está
reservada para leitura humana, para medir quantas confirmações erradas
passaram. Os números da Secção~\ref{sec:corpus-b} usam as correções aceites e
o valor original nos itens pendentes; a regeneração depois das decisões é
automática (caixa C3-06).

\paragraph{Limites.} O verificador pertence à mesma família de modelos do
extrator, e erros partilhados podem passar sem sinal. A verificação cobre
valores com citação; uma ausência (``não reportado'') só se confirma lendo as
secções relevantes do artigo, o que não foi feito. Não há conjunto de
referência humano, pelo que não se estima um erro por campo. Chamar a esta
camada ``verificada'' seria impreciso: é uma releitura independente do
extrator, com adjudicação por regra e decisão humana nos casos que a regra
não resolve.

\section{O Corpus~B, estudo a estudo}
\label{sec:apB-tabela}

A Tabela~\ref{tab:apB-corpus-b} dá, para cada um dos 39 registos, o nível mais
forte declarado nas quatro condições da Figura~\ref{fig:mapa-evidencia} e o
número de mecanismos e adaptações. É a versão em tabela do mapa de evidência.

%% Gerado por stage_c/figuras_cap3/tabela_apendice_b.py (2026-09-17). Nao editar a mao.
{\footnotesize
\begin{xltabular}{\linewidth}{@{}>{\raggedright\arraybackslash}X>{\raggedright\arraybackslash}p{2.3cm}rllllrr@{}}
\caption[Corpus B: características por estudo]{Corpus~B, estudo a estudo (39 registos cartografados). Nível mais forte entre os modelos de cada estudo. Fronteira (C5): exata, alocada, desvio reconhecido ou não reportada; re-estimação (C6): regime declarado ou não; domínio (C7); fora da amostra (C8): validação com resultado reportado. Mec.\ e Adapt.: número de ocorrências. Valores com as decisões de verificação de 2026-09-17 aplicadas e os itens pendentes no valor original.}\label{tab:apB-corpus-b}\\
\toprule
Estudo & Setor & Mod. & Fronteira & Re-est. & Domínio & Fora am. & Mec. & Adapt. \\
\midrule
\endfirsthead
\toprule
Estudo & Setor & Mod. & Fronteira & Re-est. & Domínio & Fora am. & Mec. & Adapt. \\
\midrule
\endhead
\bottomrule
\endfoot
\textcite{Alcantara2018} & cimento e cal & 4 & exata & não & não & n.a. & 0 & 0 \\
\textcite{Almanza2025} & refinação & 1 & exata & não & não & reportada & 4 & 1 \\
\textcite{Andersson2021} & pasta e papel & 2 & alocada & não & não & n.a. & 4 & 3 \\
\textcite{Bassetti2025} & gases industriais & 4 & exata & declarado & numérico & não & 5 & 3 \\
\textcite{Beisheim2019} & petroquímica & 3 & exata & declarado & qualitativo & não & 7 & 5 \\
\textcite{Beisheim2020} & petroquímica & 4 & exata & declarado & qualitativo & não & 6 & 3 \\
\textcite{Benmakhlouf2025} & química contínua & 2 & n.r. & não & não & não & 1 & 1 \\
\textcite{Bonfa2017} & vários & 4 & desvio & não & qualitativo & não & 8 & 5 \\
\textcite{Boyd2013} & cimento e cal & 3 & n.r. & declarado & não & não & 5 & 4 \\
\textcite{Boyd2016} & siderurgia e metais & 2 & alocada & não & não & não & 5 & 3 \\
\textcite{Boyd2017} & vários & 3 & exata & não & não & não & 7 & 5 \\
\textcite{Bruni2021} & cimento e cal & 11 & alocada & não & numérico & não & 6 & 4 \\
\textcite{Cagman2022} & vários & 5 & n.r. & não & qualitativo & não & 2 & 1 \\
\textcite{CastrillonMendoza2013} & cimento e cal & 4 & n.r. & não & não & não & 3 & 2 \\
\textcite{ElMaghraby2024} & química contínua & 6 & exata & não & numérico & reportada & 0 & 2 \\
\textcite{EspinelBlanco2020} & pasta e papel & 6 & exata & não & não & não & 4 & 0 \\
\textcite{Giacone2012} & siderurgia e metais, vidro & 4 & exata & não & não & não & 11 & 6 \\
\textcite{Gonzalez2012} & cimento e cal & 4 & n.r. & não & não & não & 2 & 1 \\
\textcite{Guo2020} & siderurgia e metais & 3 & n.r. & não & não & não & 2 & 1 \\
\textcite{Hamedi2019} & petroquímica & 6 & n.r. & não & não & reportada & 10 & 5 \\
\textcite{Hashim2019} & refinação & 1 & n.r. & não & não & n.a. & 1 & 1 \\
\textcite{Hashim2022} & petroquímica, vários & 2 & n.r. & não & não & n.a. & 0 & 0 \\
\textcite{Ikeyama2017} & refinação & 1 & n.r. & declarado & não & não & 2 & 2 \\
\textcite{Khalid2024a} & siderurgia e metais & 10 & alocada & não & não & n.a. & 1 & 1 \\
\textcite{Khalid2024b} & siderurgia e metais & 2 & n.r. & não & não & não & 2 & 2 \\
\textcite{LeAnh2023} & pasta e papel & 3 & exata & não & não & n.a. & 3 & 2 \\
\textcite{Luknongbu2021} & vidro & 2 & exata & não & não & não & 2 & 2 \\
\textcite{Maggiore2020} & vários & 1 & n.r. & não & não & n.a. & 2 & 1 \\
\textcite{Mahesthi2020} & cimento e cal & 6 & n.r. & declarado & não & não & 4 & 1 \\
\textcite{Moghadasi2025} & refinação & 1 & n.r. & não & não & reportada & 1 & 2 \\
\textcite{Morales2018} & petroquímica & 5 & n.r. & não & não & não & 2 & 1 \\
\textcite{Ochoa2017} & petroquímica & 4 & n.r. & não & não & não & 3 & 3 \\
\textcite{Ottermann2011} & cimento e cal & 4 & n.r. & não & não & n.a. & 1 & 1 \\
\textcite{Pelser2018} & cimento e cal & 6 & n.r. & declarado & não & n.a. & 1 & 4 \\
\textcite{Scripcariu2021} & cimento e cal, siderurgia e metais & 4 & n.r. & não & não & n.a. & 1 & 1 \\
\textcite{Sivill2013} & vários & 1 & desvio & não & não & n.a. & 6 & 2 \\
\textcite{Tarasovskii2014} & siderurgia e metais & 4 & n.r. & declarado & não & não & 4 & 3 \\
\textcite{Velazquez2013} & refinação & 1 & n.r. & não & não & não & 1 & 2 \\
\textcite{Yao2016} & petroquímica & 1 & n.r. & declarado & qualitativo & não & 3 & 2 \\
\end{xltabular}
}


\section{O que falta para uma validação humana}
\label{sec:apB-falta}

Para que os valores do Corpus~B possam ser apresentados como verificados, faltam
três passos, por ordem de custo: decidir os treze itens e as duas regras
pendentes e ler a amostra de controlo; ler, em cada publicação, as secções em
que se apoiam as ausências usadas nos resultados; e ler às cegas as onze
publicações nucleares previstas na emenda de 2026-09-13. Para o Corpus~A,
falta a auditoria por amostragem. Nenhum destes passos muda o método do
Capítulo~\ref{cha:revisao}; mudam o nível de confiança com que os seus
números podem ser citados.
